\section{Positivities, physical spaces, and spectral transport}
Positive geometry and gauge theory receive distinct structures from
the same genealogical process. The former realizes the weights and their
incidences through ordered minors, cells, and canonical forms.
The latter preserves gauge incidence, holonomy,
representations, routes, and memory within a compatible measure.
Its form of positivity relevant to physical reconstruction is
reflection positivity. The conceptual connection acquires mathematical
content when the maps producing each form are preserved.

In the geometric part, it has been proved that subdividing a cell
preserves its canonical form through cancellation of interior poles.
For nodal energies, Schur reduction preserves an effective
form and records the detail. In the gauge part, refinement
compatibility must also preserve the vacuum, the invariant subspace,
the generator, and the spectral witness. This is the chain developed
below: arithmetic--geometric scale, common measure, semigroup,
physical projection, refinement, and continuous reconstruction.

The norm of a vector, the positivity of a moment matrix, and the
existence of a Hamiltonian gap have different domains.
A congruence with an invertible factor preserves the positivity and
inertia of a form; an isometric intertwining of the generator, together
with its domain, transports its realization onto the reduced image of the
inclusion. The uniform bound and the persistent witness then determine the
limiting gap. The development therefore keeps explicit the
inclusions and semigroup identities that justify passage to the
limit. The gap scale originates in areal--volumetric incidence
and its declared dimensional unit.

The following sections bring together the projective construction and its
subsequent development depending on the coupling \(g\). The initial
comparisons with a Gibbs parametrization are subsequently completed through
the explicit action, measure, and transport. The expository order
preserves the proofs of the process and places the complete formulation
after its dependencies. The four reflections correspond to the four
Euclidean directions of the reconstruction. The development assembled here
concerns this gauge theory and its gap, with the compact,
connected, simple group declared in the theorems.

The planar supersymmetric sector in which amplitudes are studied through
amplituhedra retains a specific physical question. Connecting it with
the present realization requires the supercharge operators, their
domains, and the corresponding representation. The dimension four of the
Euclidean chart and the difference four of the modular block remain
previously defined results of their respective operators.
